% SPDX-FileCopyrightText: 2025 IObundle
%
% SPDX-License-Identifier: MIT

The Quartus and Vivado FPGA compilation toolchains are supported via \texttt{.tcl}
scripts invoked by a Makefile. The script compiles and elaborates the design for
a given set of target FPGA boards. Additional boards can be supported by
following the flow provided.

The core is instantiated within a Verilog wrapper, which depends on the
specifics of the FPGA device and the FPGA board. Timing constraints file(s) are
provided.

After compilation, reports on FPGA resource usage, power consumption, and timing
closure are generated. A post-synthesis Verilog file is created, which can be
used in post-synthesis simulation.

% SPDX-FileCopyrightText: 2025 IObundle
%
% SPDX-License-Identifier: MIT

As with ASIC synthesis (Section~\ref{sec:synth}), the FPGA build script is
the generic Vivado flow shared by every
\href{https://github.com/IObundle/py2hwsw}{py2hwsw}-managed core,
parameterized per target board by a \texttt{board.tcl}/\texttt{board.mk}
pair (for example, \texttt{iob\_aes\_ku040\_db\_g} for a Xilinx KU040
board); adding a new board means adding a new pair, not changing the flow
itself. That flow reads the sources and board constraints, runs synthesis
followed by place and route, and reports utilization, timing, and
clock-domain-crossing results before writing out the bitstream:

\begin{lstlisting}[language=tcl,caption={py2hwsw's \texttt{hardware/fpga/vivado/build.tcl}: the Xilinx Vivado flow underlying this core's FPGA build script (condensed).},label={lst:vivado_flow}]
# Read the design's Verilog/SystemVerilog sources.
foreach file [split $VSRC \ ] {
    read_verilog -sv $file
}

# Read board-specific properties and design constraints (XDC).
source vivado/$BOARD/board.tcl
read_xdc vivado/$BOARD/$NAME\_dev.sdc

# Synthesize, then implement: optimize, place, and route.
synth_design -include_dirs ../src -part $PART -top $FPGA_TOP
opt_design
place_design
route_design -timing

# Reports: utilization, timing, and clock-domain-crossing checks.
report_utilization    -file reports/$FPGA_TOP\_$PART\_utilization.rpt
report_timing_summary -file reports/$FPGA_TOP\_$PART\_timing_summary.rpt
report_cdc -details   -file reports/$FPGA_TOP\_$PART\_cdc.rpt

# Final bitstream (or, for out-of-context synthesis, a netlist and stub).
write_bitstream -force $FPGA_TOP.bit
\end{lstlisting}

Unlike ASIC synthesis, FIFO and RAM subblocks are not black-boxed here:
Versat's own two-port and dual-port memories (Section~\ref{sec:ifsig}) map
directly onto the target FPGA's block RAMs, so the utilization report
above already accounts for them.


\subsubsection*{System-level FPGA Run}

Section~\ref{sec:simple_example_software} already shows this in practice: the \href{https://github.com/IObundle/iob-soc-versat}{iob-soc-versat} repository embeds the generated accelerator in a complete RISC-V SoC and exercises it with bare-metal C software on FPGA, through \texttt{make fpga-run TEST=EXAMPLE\_Simple}. Section~\ref{sec:py2hwsw_example} works through the accelerator standing alone on FPGA instead, without a surrounding SoC, through Py2HWSW's own \texttt{make fpga-build}.
